\begin{abstract}
去蜂窝大规模多输入多输出（CF-mMIMO）系统通过密集部署大量接入点（AP），可实现无缝通信覆盖，显著提升用户频谱效率与系统总容量。然而，该架构在实际部署中面临高昂的成本压力，且不可避免地引入非理想硬件因素。本文针对上述挑战，系统研究了非理想硬件条件下CF-mMIMO系统的资源分配与性能优化问题。

具体而言，本文首先研究了在用户设备（UE）端模拟非线性功率放大器（PA）及AP端采用低分辨率模数转换器（ADC）条件下，上行CF-mMIMO系统的可达速率性能。推导了上行用户可达速率的闭式表达式，系统分析了AP数量、UE密度、AP天线数目及ADC分辨率等关键参数对系统性能的影响。为有效抑制UE间干扰并提升系统和速率，提出了一种基于图神经网络（GNN）辅助的演员-评论家（Actor-Critic）功率分配算法（DMAGNN-AC）。该框架突破了深度强化学习在高维非结构化状态空间下的表征瓶颈，实现了具备物理可解释性的特征嵌入。与全功率发射方案相比，所提功率分配策略可使系统速率提升近一倍。进一步地，为缓解低分辨率ADC对系统速率的不利影响，设计了多智能体深度Q集成网络（MADQIN）以优化ADC分辨率分配策略。

在此基础上，本文进一步研究了非理想硬件条件下CF-mMIMO系统的超可靠低时延通信（URLLC）服务资源分配策略。针对工业自动化等物联网关键应用对毫秒级时延与极高可靠性的严苛需求，构建了非线性PA的损伤模型，推导了有限块长下上行链路UE速率闭合表达式。为降低传统优化算法的计算复杂度，提出基于GNN与深度强化学习融合的智能功率控制策略，通过图嵌入技术建模UE空间相关性，动态抑制硬件损伤与时变信道干扰。所提方案在保持毫秒级时延的同时，为工业自动化等物联网应用提供理论与技术支撑。仿真结果充分验证了所提方案的有效性。

\keywords{去蜂窝大规模MIMO; 非理想硬件; 图神经网络; 深度强化学习; 功率分配; 超可靠低时延通信; 有限块长; 资源分配}
%({\color{blue}{一般选3 \textasciitilde 8个单词或专业术语，且中英文关键词必须对应。}})
\end{abstract}

\begin{englishabstract}
Cell-free massive multiple-input multiple-output (CF-mMIMO) systems, characterized by dense deployment of access points (APs), provide seamless communication coverage with significant enhancements in spectral efficiency and overall system capacity. However, the practical implementation of this architecture faces substantial cost challenges and inevitably introduces non-ideal hardware impairments. This dissertation systematically addresses resource allocation and performance optimization challenges in CF-mMIMO systems operating under non-ideal hardware conditions.

This research first investigates the achievable rate performance of uplink CF-mMIMO systems incorporating nonlinear power amplifiers (PAs) at user equipment (UE) terminals and low-resolution analog-to-digital converters (ADCs) at APs. Closed-form expressions for uplink user achievable rates are derived, enabling comprehensive analysis of key system parameters including AP deployment density, UE distribution characteristics, antenna array configurations, and ADC quantization precision. To effectively mitigate inter-user interference and enhance system throughput, a graph neural network (GNN)-assisted actor-critic power allocation algorithm (DMAGNN-AC) is proposed. This framework overcomes the representation limitations inherent in deep reinforcement learning (DRL) when dealing with high-dimensional unstructured channel state information, achieving physically interpretable feature learning. Compared with conventional full-power transmission strategies, the proposed power allocation approach achieves approximately twofold improvement in system rate performance. Furthermore, to alleviate the performance degradation caused by low-resolution ADCs, a multi-agent deep Q-integration network (MADQIN) is designed for optimizing ADC resolution allocation strategies.

Building upon these contributions, this research extends to resource allocation strategies for ultra-reliable low-latency communication (URLLC) services in CF-mMIMO systems with hardware impairments. Addressing the stringent requirements of millisecond-level latency and ultra-high reliability in critical IoT applications such as industrial automation, impairment models for nonlinear PAs are established, and closed-form rate expressions for uplink UEs under finite blocklength constraints are derived. To reduce the computational complexity of traditional optimization algorithms, an intelligent power control strategy integrating GNN with DRL is developed. By leveraging graph embedding techniques to model spatial correlations among UEs, the proposed approach enables dynamic suppression of hardware-induced distortions and time-varying channel interference. The developed solution maintains millisecond-level latency while providing theoretical and technical support for IoT applications in industrial automation. Simulation results comprehensively validate the effectiveness of the proposed methodologies.

\englishkeywords{CF-mMIMO, non-ideal hardware, graph neural network, deep reinforcement learning, power allocation, ultra-reliable low-latency communication, finite blocklength, resource allocation}

\end{englishabstract}
